\appendix

\section{Experimental details and reproducibility}
\label{app:setup}
\label{app:experimental}

\subsection{Datasets and configurations}

The experiments use \texttt{dianping}, \texttt{wikirevision},
\texttt{weibo\_daily}, and \texttt{weibo\_tech} from GDGB. These are the four
datasets used for the temporal analysis in the main paper. All R2O statistics
are computed from event timestamps after the same preprocessing used by the
generation pipeline. Dataset access and preprocessing follow the GDGB release.

The primary analysis contains 16 preregistered configurations. Twelve use
Qwen3-30B-A3B and cover the four datasets under L0, L1, and L2. L0 reproduces
GAG General, L1 adds tool use, and L2 additionally includes memory. Four
archived GPT-4o-Mini configurations provide matched backbone comparisons on
\texttt{dianping} L0, \texttt{wikirevision} L0 and L1, and
\texttt{weibo\_tech} L0. A separate exploratory analysis contains eight
Qwen3-8B configurations covering L0 and L1 on all four datasets. Exploratory
results are excluded from statements that refer to the 16 preregistered
configurations.

\subsection{Generation protocol}

Each run conditions on a real prefix containing 1,000 edges and generates a
continuation of approximately 9,000 events in rounds of 50 events. The main
evaluation uses the corresponding real continuation from the inherited
10,000-edge evaluation window. Reference statistics are computed from the
matched real slots of the same window.

Several archived runs terminate before the nominal horizon. The archived
Qwen3-30B-A3B runs for \texttt{weibo\_tech} L1,
\texttt{weibo\_tech} L2, and \texttt{weibo\_daily} L2 are shorter than the
nominal run. The \texttt{weibo\_daily} L2 structural evaluation is excluded
because the archived structural output is contaminated, while its timing
statistics use the archived run produced before contamination. The
GPT-4o-Mini run for \texttt{wikirevision} L1 contains 7,232 events and the
GPT-4o-Mini run for \texttt{weibo\_tech} L0 contains 200 generated events.
Its timing statistics use the same matched-slot procedure at the available
horizon, while structural statistics are unavailable because the structural
evaluation lane was not run for such a short stream. Configuration-level
numbers of scored events are reported in Table~\ref{tab:artifact-controls}.

The exploratory Qwen3-8B L1 runs were stopped at round boundaries because of
generation time. Their lengths range from 4,900 to 5,696 events. For
comparisons between exploratory L0 and L1, we therefore also report a
matched horizon analysis in which the corresponding L0 stream is truncated to
the L1 event count and rescored using the same procedure.

\subsection{Backbones, serving, and efficiency}

Qwen3-30B-A3B (checkpoint \texttt{Qwen3-30B-A3B-Instruct-2507-FP8}) is served
locally with vLLM 0.20.0 (FP8 weights, tensor parallel size 1, that is, one
GPU per server instance, 32,768-token context limit) on a server with eight
NVIDIA RTX 6000 Ada (48\,GB) GPUs, 64 CPU cores, and 710\,GB RAM, which also
ran the agent harness. GPT-4o-Mini (\texttt{gpt-4o-mini-2024-07-18}, version pinned) is
accessed through a commercial API gateway; its harness ran on a 128-core,
541\,GB server whose two NVIDIA RTX A6000 (48\,GB) GPUs serve only the
retrieval embedding model. The provider does not expose its serving stack or
a determinism guarantee, so API runs are not used to estimate variance for the
local model. The exploratory Qwen3-8B runs use the same API gateway with
provider reasoning disabled. All three backbones use the same decoding
configuration in the generation harness (recorded in the harness model
configuration of every cell): temperature 0.3 and at most 1,024 generated
tokens per call. Wall-clock generation times were not retained in
the run archive; per-call token and call counts are retained for the
exploratory grid (below). All scoring and analysis is CPU only.

For the exploratory Qwen3-8B grid, API accounting gives 1.9 to 2.8 LM calls per
event for L0 and 2.5 to 4.7 calls per event for L1. L1 uses 1.6 to 2.4$\times$
more prompt tokens per event than L0. The resulting cost is
\$4.48 to \$7.77 per cell and \$46.78 for the complete exploratory grid.
Code, score tables, and the verification gate are available at
\url{https://github.com/parhamouni/r2o-process-fidelity}. Per-call token usage and call counts are retained in the supplementary
material.

\subsection{Timing baselines}

Three timing baselines use only the 1,000-edge seed prefix. The Poisson
baseline estimates a constant event rate from the seed. The empirical IET
bootstrap samples seed inter-event times independently with replacement. The
Hawkes baseline fits an exponential Hawkes process to the seed by maximum
likelihood. These baselines generate event times only and do not generate graph
structure or text. Reported baseline means and standard deviations use 20
simulation seeds.

\subsection{Registration and reproducibility}
\label{app:repro}

The primary configuration set and the main R2O analyses were preregistered.
The Qwen3-8B arm was added after registration and is labeled exploratory
throughout the paper.

All analysis code, per configuration scores, matched real-window construction,
the scheduler and count-conditioning controls, figure and table generation
scripts, and a verification gate that recomputes every statistic reported in
the paper from the committed score files are provided in the
supplementary material and in the public repository. The repository data are read only under the analysis
scripts, and every figure and table is regenerable from the committed score
files. Dataset access follows the GDGB release.

The preregistration includes hypotheses, falsifiers, exclusions, a power plan,
and stop rules, together with the corresponding deviation log. The eight
exploratory Qwen3-8B cells were generated with the same harness, with provider
reasoning disabled, temperature 0.3, and \texttt{max\_tokens}=1024. Per-call
token usage was logged, and the per event accounting above is recomputed from
those logs by the verification gate.


\section{Evaluation landscape}
\label{app:landscape}

Table~\ref{tab:evaluation-landscape} summarizes how prior graph generation,
temporal graph generation, language-model-based graph generation, and
language-model-based social simulation studies evaluate generated data. The
table is included to place the R2O requirements in the broader evaluation
landscape. It distinguishes structural, textual, marginal-temporal,
temporal-dependence, likelihood-based, and closed loop evaluation. The main
paper uses this comparison for positioning and does not rely on it for the
formal results.

\begin{table}[t]
\centering
\caption{Evaluation practices for graph simulators and LM social simulators.
Metric classes are S for snapshot structural, T for text, MT for marginal
temporal, D for temporal dependence, L for predictive or likelihood based,
and CL for closed loop rollout. Dep. records whether temporal dependence of
the generated stream is tested. CL records whether a closed loop test is
performed. Every row is explicitly cited.}
\label{tab:evaluation-landscape}
\scriptsize
\resizebox{\linewidth}{!}{%
\begin{tabular}{lllll}
\toprule
Reference & System / topic & Class & Dep. & CL \\
\midrule

\multicolumn{5}{l}{\textit{Static graph generation}} \\

\citet{you2018graphrnn}
    & GraphRNN, static autoregressive graph generator & S/L & no & no \\

\citet{bojchevski2018netgan}
    & NetGAN, random walk static graph generator & S/L & no & no \\

\citet{liao2019gran}
    & GRAN, static graph generator & S & no & no \\

\citet{chanpuriya2021edge}
    & Theory of edge independent graph models & S & no & no \\

\citet{thompson2022evaluation}
    & Critique of graph generative evaluation & S & no & no \\

\citet{obray2022evaluation}
    & Critique of graph generative evaluation & S & no & no \\

\citet{shirzad2022contrastive}
    & Graph generative evaluation & S & no & no \\


\midrule
\multicolumn{5}{l}{\textit{Temporal graph generation and benchmarks}} \\

\citet{holme2012temporal}
    & Review of temporal networks & MT/D & yes & no \\

\citet{sizemore2018dynamic}
    & Dynamic graph metrics tutorial & D & partial & no \\

\citet{zhou2020taggen}
    & TagGen, temporal interaction network generator & S & no & no \\

\citet{zeno2021dymond}
    & DYMOND, motif based dynamic network generator & S/MT & partial & no \\

\citet{zhang2021tggan}
    & TG GAN, continuous time temporal graph GAN & S/MT & no & no \\

\citet{poursafaei2022better}
    & Evaluation critique for dynamic link prediction & L & partial & no \\

\citet{gupta2022survey}
    & Survey of temporal graph representation learning and generation
    & S/L & no & no \\

\citet{clarkson2022damnets}
    & DAMNETS, Markovian network time series generator & S & no & no \\

\citet{gauvin2022randomized}
    & Randomized reference models for temporal networks & D & yes & no \\

\citet{gupta2022tigger}
    & TIGGER & S/MT & no & no \\

\citet{huang2023tgb}
    & TGB benchmark, temporal graph learning & L & no & no \\

\citet{zhang2024dtgb}
    & DTGB benchmark, DyTAG discriminative tasks & T/L & no & no \\

\citet{xiang2025efficient}
    & Learning based temporal graph simulator & S & no & no \\

\citet{du2025graphmaster}
    & GraphMaster, LM agent graph synthesis & S/T & no & no \\

\citet{chang2025structurally}
    & LM generated social networks & S & no & no \\

\citet{hosseini2025quality}
    & Quality measures for dynamic graph generative models
    & S & partial & no \\


\midrule
\multicolumn{5}{l}{\textit{LM generated graphs}} \\

\citet{ji2025gag}
    & GAG, LM multi agent DyTAG generator & S/T/MT & no & no \\

\citet{peng2025gdgb}
    & GDGB benchmark & S/T/MT & no & no \\

\citet{ji2026graphia}
    & GRAPHIA & S/T & no & no \\


\midrule
\multicolumn{5}{l}{\textit{LM social simulation}} \\

\citet{park2022simulacra}
    & Social Simulacra, LM populated prototypes & -- & no & no \\

\citet{park2023generative}
    & Generative Agents, Smallville & -- & no & no \\

\citet{gao2023s3}
    & S3, LM social network simulation & T/MT & no & no \\

\citet{tornberg2023simulating}
    & LM social media simulation for feed algorithms & -- & no & no \\

\citet{gao2024survey}
    & Survey of LM empowered agent based modeling & -- & no & no \\

\citet{rossetti2024ysocial}
    & Y Social, LM social media digital twin & S/T & no & no \\

\citet{mou2024hisim}
    & HiSim, social movement simulation & T/MT & no & no \\

\citet{yang2024oasis}
    & OASIS, million agent social media simulator & T/MT & no & no \\

\citet{chuang2024opinion}
    & LM opinion dynamics & MT & no & no \\

\citet{piao2025agentsociety}
    & AgentSociety, large scale LM agent society & -- & no & no \\

\citet{wang2025recagent}
    & RecAgent, LM user behavior simulator & -- & no & no \\

\citet{larooij2025critical}
    & Critical review of generative social simulations & -- & no & no \\

\citet{saeedi2025mosaic}
    & MOSAIC, content dissemination simulation & S/MT & no & no \\

\citet{zhang2025socioverse}
    & SocioVerse, world model with a large user pool & -- & no & no \\

\citet{tomasevic2026operational}
    & Operational validation on Y Social and Voat & S/T/MT & no & no \\

\citet{koley2025salm}
    & SALM & S/T & no & no \\

\citet{zhang2026posim}
    & POSIM, Hawkes driven public opinion simulator & T/MT & partial & no \\

\citet{mou2026society}
    & Survey of LM agent social simulation & -- & no & no \\

\citet{schwager2026operational}
    & Conditioned Comment Prediction, operational validity & T & no & no \\

\citet{deng2026hawkesllm}
    & HawkesLLM, Hawkes and LM cascade simulator & T & no & partial \\

\citet{miao2026email}
    & LM multi agent email networks with Hawkes guidance
    & MT/D/CL & partial & yes \\

\citet{zhang2026crisis}
    & Hawkes gated crisis response simulator & MT/D & yes & no \\

\bottomrule
\end{tabular}}
\end{table}


\section{R2O definitions, proofs, and synthetic validation}
\label{app:proofs}
\label{app:r2o-definitions}

\subsection{Metric definitions}

For an actor with IET sequence $\tau_1,\ldots,\tau_k$, burstiness is
\[
B=\frac{\sigma_\tau-\mu_\tau}{\sigma_\tau+\mu_\tau},
\]
where $\mu_\tau$ and $\sigma_\tau$ are the mean and standard deviation of the
IETs.
This definition follows the standard burstiness formulation \citep{goh2008burstiness}. The memory coefficient is the lag 1 Pearson correlation between
consecutive IETs, computed within actors and pooled using the available
adjacent pairs. The use of consecutive intervals follows prior work on temporal memory \citep{goh2008burstiness,karsai2012universal}. Global lag 1 autocorrelation applies the same dependence
calculation to the globally ordered IET sequence.

The Rhythm axis uses the KS distance between pooled generated and real IET
distributions, the Wasserstein distance between log IET distributions, the
absolute burstiness difference, and node-rhythm divergence. Node-rhythm
divergence is the median actor specific KS distance between the generated and
real IET distributions for the same actor. Rate uses the absolute log-duration
ratio. Order uses the absolute memory difference and the absolute difference
in global lag 1 autocorrelation.

These statistics keep the three R2O axes distinct. Rate responds directly to a
change in time scale. Rhythm statistics depend on the IET distribution but not
on the sequential order of the intervals. Order statistics depend on
consecutive intervals and are invariant to uniform rescaling of time.

\subsection{Detailed non identification construction and proof}

The compact statement in Proposition~\ref{prop:nonident} combines the two
order sensitivity results used in the main paper. The first result uses the
following explicit construction. Let $F$ be any IET distribution with finite
variance. Let $P_0$ generate each actor's IETs independently from $F$, and let
$P_\rho$ generate them from a stationary Gaussian copula AR(1) process with the
same marginal $F$ and parameter $\rho\in(0,1)$. If partner and text choices are
identical and independent of event timing, then $P_0$ and $P_\rho$ have
identical per actor IET marginals, identical burstiness, and identical
distributions for every state statistic, while $M(P_\rho)>0=M(P_0)$.

Let the latent process satisfy
\[
z_j=\rho z_{j-1}+\sqrt{1-\rho^2}\,\epsilon_j,
\]
with standard Gaussian innovations, and define
\[
\tau_j=F^{-1}\!\left(\Phi(z_j)\right).
\]
For every value of $\rho$, the marginal distribution of $\tau_j$ is $F$.
The per-actor IET multisets, burstiness, and every statistic that depends only
on the IET marginals therefore have the same distribution under $P_0$ and
$P_\rho$. Partner and text choices are identical and independent of timing, so
the time aggregated graph and every state statistic also have the same
distribution. The map $F^{-1}\circ\Phi$ is non-decreasing. Positive dependence
in the latent AR(1) sequence therefore induces positive dependence between
consecutive IETs for $\rho>0$, while the dependence vanishes when $\rho=0$.
Hence $M(P_\rho)>0=M(P_0)$.

\subsection{Shuffle invariance}

\paragraph{Proof of Proposition~\ref{prop:shuffle}}
The shuffle operator changes the order of an actor's IETs but not the edge
multiset or the IET multiset. It therefore leaves the time aggregated graph and
every state or marginal timing statistic unchanged.

For an exchangeable sequence of $k$ IETs with sample mean $\bar{\tau}$ and
population sample variance $s^2$,
\[
\sum_{i\neq j}(\tau_i-\bar{\tau})(\tau_j-\bar{\tau})=-ks^2.
\]
A uniformly selected ordered pair of distinct positions therefore has expected
cross product $-s^2/(k-1)$. The expected lag 1 Pearson correlation after a
uniform permutation is consequently
\[
-\frac{1}{k-1}+o(1/k).
\]
Pooling adjacent pairs within actors preserves the $O(1/k_u)$ behavior stated
in Proposition~\ref{prop:shuffle}.

\subsection{Separation under rate changes}

\paragraph{Proof of Proposition~\ref{prop:separation}}
Under the transformation $\tau\mapsto a\tau$ with $a>0$, both the mean and
standard deviation of the IETs are multiplied by $a$. Their ratio in the
burstiness coefficient is therefore unchanged. Pearson correlation is
invariant to positive affine rescaling, so memory and lag 1 autocorrelation are
also unchanged.

The KS distance compares $F_{a\tau}(x)=F_\tau(x/a)$ with the reference
distribution and is generally changed by $a$. The same argument applies to
actor specific KS distances. For the log transformed Wasserstein distance,
\[
\log(a\tau)=\log\tau+\log a.
\]
Translation of one distribution by $c$ changes the Wasserstein distance by at
most $|c|$. Applying this with $c=\log a$ gives
\[
\left|
W_1(\log a\tau,\log\tau')
-
W_1(\log\tau,\log\tau')
\right|
\leq |\log a|.
\]
If the two unscaled log-IET distributions differ by a location shift $c$, then
$W_1 = |c|$ before rescaling and $|c + \log a|$ after, so the change equals
$\bigl||c + \log a| - |c|\bigr|$, which attains $|\log a|$ if and only if
$c\,(c + \log a) \ge 0$ (for example $c = 1$, $\log a = -3$ gives a change of
$1$, not $3$); in particular equality holds when $c = 0$, that is, when the
two unscaled distributions coincide.

\subsection{Finite sample calibration}

\paragraph{Proof of Proposition~\ref{prop:calibration}}
Under exchangeability, the rank of $D^{SR}$ among the $N+1$ values consisting
of $D^{SR}$ and $D^{RR}_{1:N}$ is uniform. Ties are resolved against the
simulator through the greater-than-or-equal comparison in the definition of
$q_{\mathrm{RR}}$, so $q_{\mathrm{RR}}$ is at least the upper tail rank divided
by $N+1$ and the bound holds with equality in the absence of ties:
\[
\Pr(q_{\mathrm{RR}}\leq\alpha)
\leq
\frac{\lfloor \alpha(N+1)\rfloor}{N+1}
\leq \alpha.
\]
The smallest attainable value is $1/(N+1)$, which occurs exactly when the
simulator to real discrepancy exceeds every real to real discrepancy.

\subsection{Teacher-forced evaluation}
\label{rem:openloop}

\paragraph{Remark 1.}
When observed history is supplied during generation, state and marginal timing
errors cannot accumulate in the same way as in a closed loop rollout. Such
evaluation can therefore understate errors that arise through feedback over
longer horizons. The Order statistics remain informative because the ordering
of intervals generated within an evaluated continuation is still produced by
the model. R2O itself is independent of the rollout procedure and can be
applied to either setting.

\subsection{Synthetic validation}

We construct a synthetic stream with 300 actors, lognormal IET marginals, and
Gaussian-copula AR(1) dependence chosen to produce memory $M=0.164$ and
burstiness $B=0.464$. We then apply transformations that isolate the
properties used by R2O. The shuffle preserves every marginal timing statistic
while removing memory. An independent renewal process with the same IET
marginal behaves similarly. A uniform rate change moves the Rate-sensitive
distributional statistics while leaving burstiness and memory unchanged.

\begin{table}[t]
\centering
\caption{Synthetic validation of the R2O requirements. The reference stream has memory 0.164 and burstiness 0.464.}
\label{tab:synthetic-validation}
\small
\begin{tabular}{lrrrr}
\toprule
Transformation & IET KS & $W_1(\log\tau)$ & $|\Delta B|$ & $M$ \\
\midrule
Faithful replicate & 0.004 & 0.009 & 0.042 & 0.169 \\
Actor-level shuffle & 0.000 & 0.000 & 0.000 & -0.007 \\
Independent renewal with same marginal & 0.006 & 0.015 & 0.039 & -0.009 \\
Rate multiplied by 20 & 0.686 & 2.996 & 0.000 & 0.164 \\
Rate multiplied by 20 with shuffle & 0.686 & 2.996 & 0.000 & -0.022 \\
\bottomrule
\end{tabular}
\end{table}


\subsection{Control identifiability}

\paragraph{Proposition 5 (Control identifiability).}
Let $Q$ denote a deterministic or randomized measurement transformation
applied by the generation pipeline, such as timestamp quantization or
duplicate handling. If a generated process $G_\theta$ has the same
distribution as the real process $E_{\mathrm{real}}$, then $Q(G_\theta)$ and
$Q(E_{\mathrm{real}})$ also have the same distribution. A discrepancy that
remains after applying the same transformation to both streams cannot be
attributed to $Q$ alone. The same logic applies to count conditioning when
generated and real actor event counts are matched before scoring.

\paragraph{Proof.}
Equality in distribution is preserved under any measurable transformation.
Applying the same transformation to both processes therefore preserves
equality. The stated interpretation follows from the contrapositive.


\section{R2O calibration and robustness checks}
\label{app:instrument}
\label{app:robustness}

\subsection{Matched real-window reference}

For each dataset, the real reference distribution is constructed from pairs of
10,000-edge windows. Candidate pairs are matched on window duration, actor
count, edges per actor, and median positive IET. We use disjoint windows where
possible and cap reuse at three pairs for any window. Up to 60 matched pairs
are selected per dataset. When the available stream provides fewer disjoint
matches, the empirical reference has correspondingly coarser resolution.

For each R2O discrepancy $D$, the main analysis records both
\[
Z_D=
\frac{D^{SR}-\mu_D^{RR}}{\sigma_D^{RR}}
\]
and the empirical tail position $q_{\mathrm{RR}}$ defined in the main paper.
The standardized score displays the size and direction of a discrepancy in
units of real variation. The empirical tail position is the primary rank-based
description. Because matched pairs can share windows, $q_{\mathrm{RR}}$ is
interpreted descriptively in the empirical analysis. Proposition 4 gives its
idealized finite sample interpretation under exchangeability.

\begin{table}[h]
\caption{Matched real-window reference distributions. Windows contain 10,000 edges. Pairs are matched on the covariates described above, with window reuse capped at three pairs. Values report mean and standard deviation of each real to real discrepancy.}
\label{tab:rr-reference}
\centering\small
\begin{tabular}{lrrrrr}
\toprule
Dataset & Windows & Pairs & KS ($\mu\pm\sigma$) & $|\Delta M|$ ($\mu\pm\sigma$) & Rhythm ($\mu\pm\sigma$) \\
\midrule
dianping & 199 & 60 & $0.043\pm0.024$ & $0.027\pm0.026$ & $0.41\pm0.11$ \\
wikirevision & 277 & 60 & $0.036\pm0.019$ & $0.023\pm0.015$ & $0.36\pm0.03$ \\
weibo\_daily & 35 & 52 & $0.073\pm0.038$ & $0.035\pm0.029$ & $0.29\pm0.02$ \\
weibo\_tech & 10 & 13 & $0.075\pm0.036$ & $0.058\pm0.046$ & $0.32\pm0.03$ \\
\bottomrule
\end{tabular}
\end{table}


\subsection{Timestamp ties}

Round-based generation assigns the same timestamp to multiple events and can
therefore produce zero IETs. We retain the original stream and evaluate two
additional representations. Moment statistics such as memory use a
deduplicated representation in which repeated actor-and-timestamp events are
removed. Distributional statistics use a jittered representation in which ties
are spread within the corresponding generation interval while preserving the
order of distinct timestamps.

Across the primary Qwen3-30B-A3B configurations, removing duplicate
timestamps explains a median 41\% of the absolute burstiness discrepancy. A
residual with the same sign remains in all 12 configurations. On two
\texttt{dianping} configurations with small original burstiness discrepancies,
deduplication increases the discrepancy, indicating that the effect of ties is
not uniformly directional.

\subsection{Scheduler control}

The scheduler control applies the same round quantization used by the
generation pipeline to the real stream before computing Order. This produces a
real reference subject to the same timestamp resolution as the generated
stream. Quantization changes real memory from a 39\% increase to a 49\%
reduction, but the resulting real memory remains above the generated
Qwen3-30B-A3B range on all four datasets.

\begin{table}[h]
\caption{Scheduler control for the memory coefficient. Real streams are passed through the same round quantization used by the generation pipeline before memory is recomputed. $M^{\mathrm{real}}$ is computed on the full real reference window; Table~\ref{tab:complete-configs} reports real $M$ on each cell's matched reference slots and Table~\ref{tab:matched-horizon} at the matched horizon, so the real reference differs slightly across tables.}
\label{tab:scheduler-control}
\centering\small
\begin{tabular}{lrrrr}
\toprule
Dataset & $M^{\mathrm{real}}$ & $M^{\mathrm{real}}_{\mathrm{quant}}$ & Reduction & Generated $M$ (min..max) \\
\midrule
dianping & 0.102 & 0.052 & $+49$\% & $-0.276$\,..\,$-0.015$ \\
wikirevision & 0.184 & 0.144 & $+22$\% & $-0.011$\,..\,$+0.030$ \\
weibo\_daily & 0.159 & 0.221 & $-39$\% & $-0.032$\,..\,$-0.027$ \\
weibo\_tech & 0.238 & 0.244 & $-3$\% & $-0.014$\,..\,$-0.009$ \\
\bottomrule
\end{tabular}
\end{table}


\subsection{Actor-count conditioning}

Differences in actor activity may alter marginal IET comparisons even when the
timing distribution conditional on activity is similar. We therefore recompute
actor level KS after matching generated and real event counts for each actor.
The procedure uses 20 subsampling replicates. The median ratio between
count conditioned and original KS is 1.006, with values ranging from 0.886 to
1.039 across the primary configurations. The marginal discrepancy therefore
remains after event counts are matched.

\subsection{Evaluation-window sensitivity}

The inherited evaluation uses the first 10,000 edges of each dataset. For
\texttt{wikirevision}, this window is unusually early relative to the other
candidate windows. We therefore repeat the GPT-4o-Mini L0 and L1 comparison on
a representative window selected using the same matching covariates as the real
reference construction. Marginal KS decreases substantially and degree MMD
also improves, but the Order discrepancy remains. Real memory on the
representative window is 0.15, while the generated values are 0.02 and -0.06
for L0 and L1 respectively.

\begin{table}[h]
\caption{Evaluation-window sensitivity on \texttt{wikirevision} using GPT-4o-Mini. The representative and inherited initial windows use the same generation and scoring procedures. Inter-event KS here is the jitter arm statistic (timestamp ties spread within the round quantum) and therefore differs from the as scored KS of Table~\ref{tab:complete-configs} for the same runs.}
\label{tab:window-sensitivity}
\centering\small
\begin{tabular}{lrrrr}
\toprule
 & \multicolumn{2}{c}{Representative window} & \multicolumn{2}{c}{Cold-start window} \\
\cmidrule(lr){2-3}\cmidrule(lr){4-5}
Metric & L0 & L1 & L0 & L1 \\
\midrule
Inter-event KS (jitter) & 0.090 & 0.088 & 0.310 & 0.383 \\
$\log(T_g/T_r)$ & 3.83 & 4.44 & 3.98 & 3.62 \\
Degree MMD (final) & 0.078 & 0.061 & 0.151 & 0.121 \\
$M$ (dedup) & $+0.02$ & $-0.06$ & $+0.43$ & $-0.03$ \\
Zero IET fraction & 0.37 & 0.30 & 0.64 & 0.51 \\
\bottomrule
\end{tabular}
\end{table}


On this representative window, the marginal jitter-arm KS is 0.090 versus
0.310 on the inherited cold-start window, a roughly 3.4$\times$
window-specific component of marginal error. Structure is also better, while
the dependence deficit persists. Thus the inherited first-10k-edge convention
simultaneously understates structural fidelity and overstates marginal timing
error.

A particularly transparent Rate example appears on this same representative
window: the tool loop worsens the log-duration ratio from 3.83 to 4.44
(approximately $46\times$ to $85\times$ in duration) while structure improves.
This illustrates directly that Rate and structural fidelity can move in
opposite directions.

\subsection{Summary of robustness checks}

The controls address distinct alternative explanations. Timestamp handling
changes the magnitude of some Rhythm statistics but does not restore Order.
Applying the generation scheduler to real data does not account for the
missing memory. Matching actor event counts leaves the marginal discrepancy
nearly unchanged. Selecting a more representative window reduces marginal
error but leaves the dependence discrepancy. These results support the main
conclusion that the R2O findings are not explained by the measurement artifacts
tested here.


\section{Complete experimental results}
\label{app:tables}
\label{app:complete-results}

This appendix reports the complete configuration level results underlying the
main experiments. Table~\ref{tab:complete-configs} contains the primary and
exploratory configurations. Table~\ref{tab:matched-horizon} gives the
matched horizon analysis for the exploratory Qwen3-8B runs.
Table~\ref{tab:timing-baselines} reports the timing baselines, and
Table~\ref{tab:artifact-controls} reports configuration level quantities used
in the robustness analysis. Two KS statistics appear in these tables and are
not interchangeable: Table~\ref{tab:complete-configs} reports the pooled
actor level inter-event KS against matched reference slots (the R2O Rhythm
statistic), whereas Table~\ref{tab:timing-baselines} reports a global stream
inter-event KS, because the timing baselines generate a single global stream
without actor assignment. For the same run these two values differ (for
example \texttt{dianping} L0: 0.372 versus 0.823). Table~\ref{tab:artifact-controls}
additionally reports the rate normalized Rhythm statistic of
Proposition~\ref{prop:separation}: per actor inter-event times are divided by
the actor's mean interval before pooling (shape KS), which removes the Rate
component. Shape KS ranges from 0.22 to 0.66 across the 16 preregistered
configurations and exceeds the matched slot KS on 11 of them; on the five
where it is lower (all four \texttt{dianping} configurations and
\texttt{wikirevision} GPT-4o-Mini L0) it remains at or above 0.22, so the
Rhythm discrepancy is not an artifact of the Rate mismatch.

\begin{table}[t]
\caption{Complete results for the 16 preregistered configurations and eight exploratory Qwen3-8B configurations. Degree MMD reports structural error; KS is pooled actor level inter-event KS against matched reference slots (not the global stream KS of Table~\ref{tab:timing-baselines}); $Z_{\mathrm{KS}}$ is the same discrepancy in real reference units; $M$ is the memory coefficient from the deduplicated representation, with the real value in parentheses; $q_{\mathrm{RR}}$ is the empirical real reference tail position. Superscript $s$ marks a short run; $c$ marks a structural evaluation affected by contamination (\texttt{weibo\_daily} L2 excluded; \texttt{weibo\_tech} L2 scored on the archived 8.65k-edge run); $t$ marks the 200 event truncated archive, for which the structural evaluation is unavailable and timing uses matched slots at the available horizon; $e$ marks exploratory Qwen3-8B configurations added after registration (L1 runs stopped at round boundaries, 4,900 to 5,696 events, so 8B L0 and L1 rows are scored at different event counts; see Table~\ref{tab:matched-horizon}).}
\label{tab:complete-configs}
\centering\scriptsize
\setlength{\tabcolsep}{4pt}\renewcommand{\arraystretch}{0.86}
\begin{tabular}{lllrrrrr}
\toprule
Dataset & Backbone & Cell & Deg.\ MMD & KS & $Z_{\mathrm{KS}}$ & $M$ ($M^{\mathrm{real}}$) & $q_{\mathrm{RR}}$ \\
\midrule
dianping & Qwen3-30B-A3B & L0 & 0.052 & 0.372 & $+14.0$ & $-0.01$ (0.10) & 0.016 \\
 &  & L1 & 0.026 & 0.440 & $+16.9$ & $-0.26$ (0.10) & 0.016 \\
 &  & L2 & 0.034 & 0.433 & $+16.6$ & $-0.28$ (0.10) & 0.016 \\
 & GPT-4o-Mini & L0 & 0.056 & 0.340 & $+12.6$ & $-0.10$ (0.10) & 0.016 \\
 & Qwen3-8B & L0$^{e}$ & 0.059 & 0.444 & $+17.0$ & $-0.10$ (0.10) & 0.016 \\
 &  & L1$^{e,s}$ & 0.045 & 0.578 & $+22.7$ & $-0.06$ (0.11) & 0.016 \\
wikirevision & Qwen3-30B-A3B & L0 & 0.159 & 0.456 & $+22.5$ & $+0.03$ (0.20) & 0.016 \\
 &  & L1 & 0.139 & 0.283 & $+13.2$ & $+0.00$ (0.20) & 0.016 \\
 &  & L2 & 0.115 & 0.271 & $+12.6$ & $-0.01$ (0.20) & 0.016 \\
 & GPT-4o-Mini & L0 & 0.151 & 0.328 & $+15.7$ & $+0.43$ (0.20) & 0.016 \\
 &  & L1$^{s}$ & 0.121 & 0.228 & $+10.3$ & $-0.03$ (0.22) & 0.016 \\
 & Qwen3-8B & L0$^{e}$ & 0.131 & 0.580 & $+29.1$ & $+0.16$ (0.20) & 0.016 \\
 &  & L1$^{e,s}$ & 0.077 & 0.257 & $+11.8$ & $+0.01$ (0.25) & 0.016 \\
weibo\_daily & Qwen3-30B-A3B & L0 & 0.273 & 0.305 & $+6.0$ & $-0.03$ (0.16) & 0.019 \\
 &  & L1 & 0.112 & 0.339 & $+6.9$ & $-0.03$ (0.16) & 0.019 \\
 &  & L2$^{s}$ & --$^{c}$ & 0.300 & $+5.9$ & $-0.03$ (0.17) & 0.019 \\
 & Qwen3-8B & L0$^{e}$ & 0.239 & 0.299 & $+5.9$ & $+0.00$ (0.16) & 0.019 \\
 &  & L1$^{e,s}$ & 0.137 & 0.296 & $+5.8$ & $-0.04$ (0.13) & 0.019 \\
weibo\_tech & Qwen3-30B-A3B & L0 & 0.209 & 0.297 & $+6.2$ & $-0.01$ (0.22) & 0.071 \\
 &  & L1$^{s}$ & 0.097 & 0.264 & $+5.3$ & $-0.01$ (0.27) & 0.071 \\
 &  & L2$^{s}$ & 0.214$^{c}$ & 0.265 & $+5.3$ & $-0.01$ (0.27) & 0.071 \\
 & GPT-4o-Mini & L0$^{t}$ & --$^{t}$ & 0.374 & $+8.4$ & $-0.04$ (0.16) & 0.071 \\
 & Qwen3-8B & L0$^{e}$ & 0.174 & 0.273 & $+5.6$ & $+0.01$ (0.22) & 0.071 \\
 &  & L1$^{e,s}$ & 0.116 & 0.331 & $+7.2$ & $-0.01$ (0.23) & 0.071 \\
\bottomrule
\end{tabular}
\end{table}


\begin{table}[t]
\centering
\small
\caption{Matched-horizon analysis for the exploratory Qwen3-8B configurations. Each L0 stream is truncated to the event count of its paired L1 stream and rescored using the same procedure. KS uses the original generated representation and $M$ uses the deduplicated representation.}
\label{tab:matched-horizon}
\begin{tabular}{llrrrrrr}
\toprule
Dataset & Arm & $n$ & KS & $Z_{\mathrm{KS}}$ & $M_{\mathrm{gen}}$ & $M_{\mathrm{ref}}$ & $Z_{|\Delta M|}$ \\
\midrule
wikirevision & L0-full & 9000 & 0.580 & $+29.1$ & $+0.162$ & 0.203 & $+1.2$ \\
wikirevision & L0-trunc & 5696 & 0.416 & $+20.3$ & $-0.006$ & 0.246 & $+15.1$ \\
wikirevision & L1 & 5696 & 0.257 & $+11.8$ & $+0.008$ & 0.246 & $+14.2$ \\
\addlinespace
dianping & L0-full & 9000 & 0.444 & $+17.0$ & $-0.105$ & 0.102 & $+6.9$ \\
dianping & L0-trunc & 5650 & 0.494 & $+19.2$ & $-0.079$ & 0.108 & $+6.2$ \\
dianping & L1 & 5650 & 0.578 & $+22.7$ & $-0.060$ & 0.108 & $+5.4$ \\
\addlinespace
weibo\_daily & L0-full & 9000 & 0.299 & $+5.9$ & $+0.002$ & 0.157 & $+4.2$ \\
weibo\_daily & L0-trunc & 5350 & 0.348 & $+7.2$ & $-0.016$ & 0.126 & $+3.7$ \\
weibo\_daily & L1 & 5350 & 0.296 & $+5.8$ & $-0.043$ & 0.126 & $+4.7$ \\
\addlinespace
weibo\_tech & L0-full & 9000 & 0.273 & $+5.6$ & $+0.006$ & 0.225 & $+3.5$ \\
weibo\_tech & L0-trunc & 4900 & 0.275 & $+5.6$ & $-0.012$ & 0.233 & $+4.1$ \\
weibo\_tech & L1 & 4900 & 0.331 & $+7.2$ & $-0.011$ & 0.233 & $+4.1$ \\
\addlinespace
\bottomrule
\end{tabular}
\end{table}


\begin{table}[h]
\caption{Timing baseline results on the inherited evaluation window. Baselines condition only on the 1,000-edge seed prefix. Baseline values report mean and standard deviation over 20 simulation seeds. Agent configurations are single archived runs scored using the same global stream timing statistics; this global KS differs from the actor level matched slot KS of Table~\ref{tab:complete-configs}. $\rho_1$ is the global stream lag 1 autocorrelation of the generated stream, with the real value in parentheses: mean over seeds for baselines; for agent rows (Qwen3-30B-A3B) the single archived run on the deduplicated stream, with the real value at the cell's matched horizon. $\rho_1$ is a global stream statistic and need not agree with the per actor memory coefficient $M$ of Table~\ref{tab:complete-configs}. Removing tied events or collapsing the stream to distinct timestamps changes the agent values by at most 0.004 (\texttt{analysis/rho1\_tie\_check.csv}), so they do not arise from round batching.}
\label{tab:timing-baselines}
\centering\small
\begin{tabular}{llrrr}
\toprule
Dataset & Generator & Inter-event KS & Burst gap & $\rho_1$ (real $\rho_1$) \\
\midrule
dianping & seed-IET bootstrap & $0.072\pm0.004$ & $+0.025\pm0.025$ & $-0.002$ ($0.217$) \\
 & Hawkes (seed MLE) & $0.195\pm0.030$ & $-0.004\pm0.018$ & $+0.246$ ($0.217$) \\
 & Poisson (seed rate) & $0.281\pm0.005$ & $-0.510\pm0.007$ & $-0.001$ ($0.217$) \\
 & agent L0 & $0.823$ & $+0.465$ & $+0.084$ ($0.217$) \\
 & agent L1 & $0.559$ & $+0.453$ & $+0.428$ ($0.217$) \\
 & agent L2 & $0.580$ & $+0.463$ & $+0.008$ ($0.217$) \\
\addlinespace[2pt]
wikirevision & seed-IET bootstrap & $0.394\pm0.004$ & $+0.302\pm0.012$ & $-0.002$ ($0.230$) \\
 & Hawkes (seed MLE) & $0.342\pm0.039$ & $+0.327\pm0.097$ & $+0.307$ ($0.230$) \\
 & Poisson (seed rate) & $0.663\pm0.004$ & $-0.385\pm0.007$ & $-0.001$ ($0.230$) \\
 & agent L0 & $0.632$ & $+0.563$ & $+0.610$ ($0.230$) \\
 & agent L1 & $0.346$ & $+0.590$ & $+0.000$ ($0.230$) \\
 & agent L2 & $0.336$ & $+0.590$ & $+0.000$ ($0.230$) \\
\addlinespace[2pt]
weibo\_daily & seed-IET bootstrap & $0.065\pm0.004$ & $-0.377\pm0.017$ & $+0.000$ ($0.358$) \\
 & Hawkes (seed MLE) & $0.112\pm0.000$ & $-0.417\pm0.015$ & $+0.204$ ($0.358$) \\
 & Poisson (seed rate) & $0.186\pm0.005$ & $-0.684\pm0.007$ & $-0.001$ ($0.358$) \\
 & agent L0 & $0.602$ & $+0.275$ & $+0.158$ ($0.358$) \\
 & agent L1 & $0.532$ & $+0.283$ & $+0.005$ ($0.358$) \\
 & agent L2 & $0.538$ & $+0.289$ & $+0.006$ ($0.253$) \\
\addlinespace[2pt]
weibo\_tech & seed-IET bootstrap & $0.122\pm0.005$ & $-0.315\pm0.005$ & $+0.002$ ($0.301$) \\
 & Hawkes (seed MLE) & $0.172\pm0.001$ & $-0.304\pm0.009$ & $+0.144$ ($0.301$) \\
 & Poisson (seed rate) & $0.306\pm0.005$ & $-0.540\pm0.007$ & $-0.001$ ($0.301$) \\
 & agent L0 & $0.527$ & $+0.412$ & $+0.005$ ($0.301$) \\
 & agent L1 & $0.468$ & $+0.438$ & $+0.056$ ($0.297$) \\
 & agent L2 & $0.486$ & $+0.447$ & $+0.012$ ($0.296$) \\
\bottomrule
\end{tabular}
\end{table}


\begin{table}[h]
\caption{Configuration-level quantities used in the robustness analysis. $n$ is the number of generated events scored. Zero IET is the fraction of actor level inter-event times equal to zero in the generated stream. $\log(T_g/T_r)$ is the log-duration ratio. CC ratio is count conditioned KS divided by original KS. Shape KS is the pooled KS after dividing each actor's inter-event times by the actor's mean interval (rate normalized Rhythm, Proposition~\ref{prop:separation}). Rhythm is the median actor level KS against the matched real continuation.}
\label{tab:artifact-controls}
\centering\footnotesize
\begin{tabular}{lllrrrrrr}
\toprule
Dataset & Backbone & Cell & $n$ & Zero IET & $\log(T_g/T_r)$ & CC ratio & Shape KS & Rhythm \\
\midrule
dianping & Qwen3-30B-A3B & L0 & 9000 & 0.21 & $+5.11$ & 0.886 & 0.261 & 0.46 \\
 &  & L1 & 9000 & 0.16 & $+5.21$ & 0.913 & 0.408 & 0.47 \\
 &  & L2 & 9000 & 0.16 & $+5.22$ & 0.913 & 0.398 & 0.50 \\
 & GPT-4o-Mini & L0 & 9000 & 0.34 & $-0.11$ & 1.039 & 0.218 & 0.44 \\
 & Qwen3-8B & L0 & 9000 & 0.31 & $+5.23$ & 0.937 & -- & 0.56 \\
 &  & L1 & 5650 & 0.21 & $+5.67$ & 0.968 & -- & 0.61 \\
wikirevision & Qwen3-30B-A3B & L0 & 9000 & 0.67 & $+3.48$ & 1.003 & 0.491 & 0.67 \\
 &  & L1 & 9000 & 0.44 & $+3.44$ & 1.027 & 0.425 & 0.42 \\
 &  & L2 & 9000 & 0.43 & $+3.42$ & 1.025 & 0.423 & 0.43 \\
 & GPT-4o-Mini & L0 & 9000 & 0.64 & $+3.98$ & 1.003 & 0.313 & 0.42 \\
 &  & L1 & 7232 & 0.51 & $+3.62$ & 1.003 & 0.387 & 0.40 \\
 & Qwen3-8B & L0 & 9000 & 0.87 & $+4.89$ & 1.004 & -- & 0.67 \\
 &  & L1 & 5696 & 0.74 & $+4.15$ & 1.007 & -- & 0.35 \\
weibo\_daily & Qwen3-30B-A3B & L0 & 9000 & 0.36 & $+6.50$ & 0.998 & 0.662 & 0.33 \\
 &  & L1 & 9000 & 0.39 & $+6.34$ & 1.029 & 0.613 & 0.43 \\
 &  & L2 & 7500 & 0.35 & $+6.72$ & 1.010 & 0.637 & 0.38 \\
 & Qwen3-8B & L0 & 9000 & 0.29 & $+6.52$ & 0.956 & -- & 0.52 \\
 &  & L1 & 5350 & 0.34 & $+7.08$ & 1.006 & -- & 0.54 \\
weibo\_tech & Qwen3-30B-A3B & L0 & 9000 & 0.45 & $+6.44$ & 1.000 & 0.647 & 0.33 \\
 &  & L1 & 7179 & 0.42 & $+6.22$ & 1.021 & 0.591 & 0.34 \\
 &  & L2 & 7650 & 0.42 & $+6.17$ & 1.015 & 0.594 & 0.33 \\
 & GPT-4o-Mini & L0 & 200 & 0.46 & $+9.67$ & 1.016 & 0.566 & 0.62 \\
 & Qwen3-8B & L0 & 9000 & 0.42 & $+6.26$ & 1.051 & -- & 0.50 \\
 &  & L1 & 4900 & 0.49 & $+6.81$ & 1.124 & -- & 0.43 \\
\bottomrule
\end{tabular}
\end{table}



\section{Cross system scope and GRAPHIA}
\label{app:graphia}

GRAPHIA generates activity in fixed one-day bins. An activity predictor
determines counts, while the language-model component determines destinations
and text; no continuous per event timestamp exists in its generation path.
Continuous-time process fidelity is therefore structurally outside its output
space. Strong performance under current DyTAG evaluation consequently does not
require a continuous time process claim. For the same reason, GRAPHIA cannot
provide an independent-codebase replication of the residual event level timing
deficit studied here.


\section{Preregistered confirmatory tests: roll forward and timing locus}
\label{app:followups}

The experiments reported in this paper establish that the evaluated simulators
do not reproduce the real temporal process and that this failure is separable
from structural fidelity. Two confirmatory tests follow from that finding. The
\emph{future blind roll forward test} asks whether the temporal deficit grows
when the simulator is rolled forward on its own generated history instead of
being re-anchored to real history at every round. The \emph{timing locus test}
asks which component of the agent stack is responsible for the deficit, by
comparing where in the stack event timing is decided. Both protocols are frozen
and git timestamped. They are protocol only; no results from either test are
reported in this paper.

\subsection{Future blind roll forward test}

The roll forward test compares paired arms from shared warm-starts. The future blind closed loop
arm generates
\[
\widehat{G}_{t+1}
\sim
P_\theta\!\left(
\cdot \mid
G_{\le t_c},
\widehat{G}_{t_c+1:t}
\right),
\]
where the real continuation after the cutoff is used only for scoring. The
closed arm's actor schedule is recency-weighted sampling from its own generated
history and is count-matched by construction.

The teacher forced open loop arm generates
\[
\widehat{G}_{t+1}
\sim
P_\theta\!\left(\cdot\mid G_{\le t}\right),
\]
with real post-cutoff history re-injected each step. It is deliberately not
future blind.

Define
\[
\Delta(h)
=
S_{\mathrm{temporal}}^{\mathrm{closed}}(h)
-
S_{\mathrm{temporal}}^{\mathrm{open}}(h).
\]
The primary estimand is
\[
\mathrm{AUC}_{\Delta},
\]
the area under the closed-minus-open temporal-composite gap as a function of
horizon. The uncertainty interval is a paired seed level 95\% bootstrap
interval. A confidence interval containing zero is the preregistered falsifier
for the feedback explanation: it indicates temporal-process failure rather
than a failure caused by roll forward feedback.

Horizon checkpoints are equally spaced in event count, the protocol's native
clock. A growth-shape parameter is secondary. Seed noise was measured from 24
replicate sets. Paired designs with shared warm-starts are mandatory and
require approximately 4 to 5$\times$ fewer runs at $\rho=0.8$ than an unpaired
design. Full-protocol variance is measured using at least three seeds and at
least two arms before any full tier sample size is fixed.

\subsection{Preregistered temporal composites}

Four composites are preregistered:
\[
S_{\mathrm{micro}},\qquad
S_{\mathrm{temporal}},\qquad
S_{\mathrm{macro}},\qquad
S_{\mathrm{individual}}.
\]
Each is a mean of direction normalized $Z$ scores against the matched
real to real band, with larger values indicating worse mismatch.

The temporal composite contains the marginal members
$W_1(\log \mathrm{IET})$ and the absolute burstiness gap, together with a
separate preregistered temporal subblock consisting of $\Delta M$ and
node-rhythm divergence. The preregistration originally called the latter pair
the ``dependence'' block. Under the definitions used in this paper, however,
only $\Delta M$ is order-sensitive; node-rhythm divergence is a
marginal timing statistic. We therefore retain the preregistered membership
while referring to the pair as the preregistered temporal subblock.
Composites are never reported without their component panel.

The archived exploratory results use KS as the headline marginal statistic
because the archived scores predate the log-$W_1$ convention. KS remains a
preregistered diagnostic, and confirmatory runs compute both KS and log-$W_1$.

The confirmatory real to real band uses a block bootstrap preserving actor
identity, with at least 30 covariate-matched pairs where disjoint windows
suffice. \texttt{weibo\_tech} has only 13 disjoint pairs, so its $q_{\mathrm{RR}}$
floor is $1/14 = 0.071$ and every cell sits at that floor; the conclusion on
that dataset rests on the standardized $Z$ scores ($Z_{\mathrm{KS}} \ge +5.3$,
$Z_{|\Delta M|} \ge +2.5$) and on the block bootstrap band. Bootstrap replicates
estimate band uncertainty only; they do not add pseudo-observations to the
$q_{\mathrm{RR}}$ denominator.

\subsection{Timing locus test}

The timing locus test compares a predeclared menu of timing loci:
\begin{enumerate}
    \item intensity determined from history;
    \item intensity determined from state and history;
    \item direct LM emission of $\Delta t$; and
    \item a small and large handoff in which the expensive component retains
    semantic decisions about whom and what, while a lightweight temporal model
    determines when.
\end{enumerate}

The small and large factorization is not itself claimed as novel. ``Handoff wins''
is preregistered as an acceptable outcome, including the case in which direct
LM emission of $\Delta t$ loses. Non-inferiority guardrails apply to
$S_{\mathrm{macro}}$ and edge text similarity. The baseline set also includes
a trained neural temporal generator, actor level variants, and cost
normalization.

\subsection{E-REP stage 2}

The confirmatory replication of the window sensitivity result with the local
backbone is judged on a preregistered falsifier. The falsifier fires only if
both the structural/temporal decoupling and the residual timing gap vanish on
the representative window.


\section{Deviation log}
\label{app:deviations}

The preregistration carries a git timestamped deviation log. Every entry was
committed before the data it governs existed.

Entries 1 to 5 concern wording, band construction (including the requirement for
a block bootstrap where disjoint real-window pairs are insufficient),
$Z$-scoring of $W_1$ on log IETs, the two-stage E-REP split, and an E-REP
second dataset rule.

Entries 6 to 8 fix the roll forward test's primary estimand
($\mathrm{AUC}_{\Delta}$ with uniform weights, directional H1, and paired
seed level bootstrap) and make the count versus timing decomposition primary.
These entries postdate the exploratory results reported in this paper but
precede any roll forward data.

Entry 7 makes $q_{\mathrm{RR}}$ the primary real reference report, with
$Z$ retained for visualization. Applied retroactively to the exploratory
results as a robustness view, this choice cannot change any conclusion because
all 16 preregistered cells are outside the real reference band on
inter-event KS under both conventions.

Entry 9 fixes the closed arm actor schedule as recency-weighted sampling from
its own generated history, count-matched by construction.

Entry 10 was added after the main results and records the eight exploratory
Qwen3-8B cells (L0/L1 on all four datasets). They were generated under the same
seed prefix, evaluation window, decoding temperature, and scoring pipeline.
They are API served with provider reasoning disabled. Retrieval embeddings are
computed on CPU and are numerically equivalent. All four L0 lanes completed the
full horizon; each L1 lane was stopped at a round boundary
(4,900 to 5,696 events) when round latency made the full horizon uneconomical.
They are labeled exploratory everywhere and are excluded from every
``16-cell'' and ``12-cell'' statement.

The full log, including dates, rationales, and affected claims, is included in
the supplementary material.


\section{Extended limitations and scope}
\label{app:limitations}

All reported results are scored under the lineage's teacher forced protocol.
Because real history is re-injected during generation, state and per-actor
marginal errors cannot compound in the same way as under a future blind closed
loop rollout. The reported Order results should therefore be read as optimistic
with respect to process error. This argument is specific to Order:
under closed loop rollout agents condition on their own generated histories,
which may be more regular than real histories, so Rate and Rhythm could improve
on some axes while Order worsens. Whether closed loop feedback produces such
trade offs is exactly what the preregistered roll forward test
(Appendix~\ref{app:followups}) measures; it has not yet been run.

The Order axis uses lag 1 statistics ($M$ and $\rho_1$). Two streams could
agree on these while differing in higher order, longer range, or nonlinear
dependence. Lag $k$ autocorrelation, mutual information between successive
intervals, or a conditional intensity fit can be added as further Order
statistics and calibrated against the same real to real reference without
changing the framework.

The empirical claims are scoped to one lineage, one generation harness, four
datasets, and 16 preregistered cells. The backbone comparison covers the three
tested backbones, which span two model families. The Qwen3-8B arm is exploratory, single-run, and its L0/L1
contrasts are also rescored at matched horizon.

Agent configurations are single archived runs rather than replicated
generation seeds. Five cells are shorter than the nominal horizon (Table~\ref{tab:artifact-controls}). API served
backbones carry provider side serving opacity. The matched real to real
calibration quantifies discrepancy relative to natural real-stream variability
but does not estimate generation seed uncertainty. The empirical
$q_{\mathrm{RR}}$ value is a tail position determined by the available
real-window reference set and should not be interpreted as a hypothesis test
$p$-value in the empirical analysis.

Cross-system generalization is additionally constrained by output-space
differences. In particular, GRAPHIA does not expose continuous per event time,
so it cannot serve as an independent event level replication of the
continuous time residual studied here. Appendix~\ref{app:graphia} gives the
full scope argument.

The preregistration and its deviations are described in
Appendix~\ref{app:deviations}.


\section{Broader impact}
\label{app:broader-impact}

This work concerns evaluation methodology. Its intended effect is protective:
agent-based graph simulators are beginning to be used for forecasting,
what if analysis, and synthetic data generation, while the results here show
that commonly used validation can miss properties on which those uses depend.
Wider adoption of process level evaluation should reduce overtrust in simulator
outputs.

A possible negative effect is the converse: better temporal-fidelity
evaluation could eventually help build simulators of human activity that are
harder to distinguish from real traces. The roll forward protocol itself is a
detection style instrument and is dual use in the same limited sense as other
fidelity metrics.
